Em resumo, a aplicação de técnicas de aprimoramento sobre modelos de LLM, como \textit{Zero-shot}, \textit{Few-shot} e RAG, mostraram-se caminhos relevantes para explorar aspectos capazes de melhorar a geração dos modelos. Isso ocorre ao definir um papel para tornar a resposta mais específica, através de uma persona de autoridade, ao estabelecer o padrão de resposta esperado, por meio de exemplos que definem a estrutura, e, por último, ao utilizar documentos do domínio do problema, como o de gado leiteiro, fazendo com que as respostas sejam baseadas em documentos que são referências no assunto, facilitando a confirmação de fatos.

Contudo, ao avaliar os resultados obtidos, o uso de \textit{Zero-shot} nos modelos de LLM mostrou-se uma técnica capaz de gerar respostas melhores em comparação às que tiveram a aplicação de sistemas RAG. Um grande fator que pode ter influenciado esse resultado inferior foi a quantidade limitada de documentos, delimitando as respostas apenas aos dados neles presentes e diminuindo sua eficácia. Adicionalmente, o avanço de modelos que implementam busca nativa, como as versões recentes do ChatGPT e Gemini, sugere que parte desses dados já possa estar integrada à base de conhecimento dos modelos. Outro fator que impactou nos resultados do RAG, especificamente no caso do \textit{Light RAG} foi a complexidade computacional para adicionar um novo documento ao sistema, como também a alta latência e custo por indexação, restringindo a avaliação a apenas 2 modelos e com um conjunto de dados bem limitado. Essas limitações indicam o ganho significativo de sistemas complexos de recuperação estão condicionadas a quantidade e qualidade da base de conhecimento fornecida e à viabilidade de recursos computacionais.

Com base nessas observações, como trabalhos futuros, sugere-se analisar como o uso de técnicas de ajuste fino (Fine-Tuning) pode impactar a geração de respostas em modelos treinados no domínio específico da pecuária leiteira, em comparação aos resultados obtidos com Zero-shot e sistemas de RAG, considerando as métricas utilizadas neste trabalho.

Além disso, sugere-se a utilização de uma maior quantidade de dados apurados por especialistas da área, como artigos, revistas e apresentações científicas sobre pecuária leiteira, avaliando o impacto que a quantidade de documentos exerce no resultado final das respostas geradas por um modelo base, a fim de identificar se a qualidade e a quantidade dos documentos coletados são relevantes para responder adequadamente às perguntas propostas.

O sistema desenvolvido neste trabalho está disponível publicamente para reprodução e extensão em: \url{https://gitlab.com/Zander404/TCCs_DairyColletor.git}
.


